%Version 3.1 December 2024
% See section 11 of the User Manual for version history
%
%%%%%%%%%%%%%%%%%%%%%%%%%%%%%%%%%%%%%%%%%%%%%%%%%%%%%%%%%%%%%%%%%%%%%%
%%                                                                 %%
%% Please do not use \input{...} to include other tex files.       %%
%% Submit your LaTeX manuscript as one .tex document.              %%
%%                                                                 %%
%% All additional figures and files should be attached             %%
%% separately and not embedded in the \TeX\ document itself.       %%
%%                                                                 %%
%%%%%%%%%%%%%%%%%%%%%%%%%%%%%%%%%%%%%%%%%%%%%%%%%%%%%%%%%%%%%%%%%%%%%

%%\documentclass[referee,sn-basic]{sn-jnl}% referee option is meant for double line spacing

%%=======================================================%%
%% to print line numbers in the margin use lineno option %%
%%=======================================================%%

%%\documentclass[lineno,pdflatex,sn-basic]{sn-jnl}% Basic Springer Nature Reference Style/Chemistry Reference Style

%%=========================================================================================%%
%% the documentclass is set to pdflatex as default. You can delete it if not appropriate.  %%
%%=========================================================================================%%

%%\documentclass[sn-basic]{sn-jnl}% Basic Springer Nature Reference Style/Chemistry Reference Style

%%Note: the following reference styles support Namedate and Numbered referencing. By default the style follows the most common style. To switch between the options you can add or remove “Numbered” in the optional parenthesis. 
%%The option is available for: sn-basic.bst, sn-chicago.bst%  
 
%%\documentclass[pdflatex,sn-nature]{sn-jnl}% Style for submissions to Nature Portfolio journals
%%\documentclass[pdflatex,sn-basic]{sn-jnl}% Basic Springer Nature Reference Style/Chemistry Reference Style
\documentclass[pdflatex,sn-mathphys-num]{sn-jnl}% Math and Physical Sciences Numbered Reference Style
%%\documentclass[pdflatex,sn-mathphys-ay]{sn-jnl}% Math and Physical Sciences Author Year Reference Style
%%\documentclass[pdflatex,sn-aps]{sn-jnl}% American Physical Society (APS) Reference Style
%%\documentclass[pdflatex,sn-vancouver-num]{sn-jnl}% Vancouver Numbered Reference Style
%%\documentclass[pdflatex,sn-vancouver-ay]{sn-jnl}% Vancouver Author Year Reference Style
%%\documentclass[pdflatex,sn-apa]{sn-jnl}% APA Reference Style
%%\documentclass[pdflatex,sn-chicago]{sn-jnl}% Chicago-based Humanities Reference Style

%%%% Standard Packages
%%<additional latex packages if required can be included here>

\usepackage{graphicx}%
\usepackage{multirow}%
\usepackage{amsmath,amssymb,amsfonts}%
\usepackage{amsthm}%
\usepackage{mathrsfs}%
\usepackage[title]{appendix}%
\usepackage{xcolor}%
\usepackage{textcomp}%
\usepackage{manyfoot}%
\usepackage{booktabs}%
\usepackage{algorithm}%
\usepackage{algorithmicx}%
\usepackage{algpseudocode}%
\usepackage{listings}%
%%%%

%%%%%=============================================================================%%%%
%%%%  Remarks: This template is provided to aid authors with the preparation
%%%%  of original research articles intended for submission to journals published 
%%%%  by Springer Nature. The guidance has been prepared in partnership with 
%%%%  production teams to conform to Springer Nature technical requirements. 
%%%%  Editorial and presentation requirements differ among journal portfolios and 
%%%%  research disciplines. You may find sections in this template are irrelevant 
%%%%  to your work and are empowered to omit any such section if allowed by the 
%%%%  journal you intend to submit to. The submission guidelines and policies 
%%%%  of the journal take precedence. A detailed User Manual is available in the 
%%%%  template package for technical guidance.
%%%%%=============================================================================%%%%

%% as per the requirement new theorem styles can be included as shown below
\theoremstyle{thmstyleone}%
\newtheorem{theorem}{Theorem}%  meant for continuous numbers
%%\newtheorem{theorem}{Theorem}[section]% meant for sectionwise numbers
%% optional argument [theorem] produces theorem numbering sequence instead of independent numbers for Proposition
\newtheorem{proposition}[theorem]{Proposition}% 
%%\newtheorem{proposition}{Proposition}% to get separate numbers for theorem and proposition etc.

\theoremstyle{thmstyletwo}%
\newtheorem{example}{Example}%
\newtheorem{remark}{Remark}%

\theoremstyle{thmstylethree}%
\newtheorem{definition}{Definition}%

\raggedbottom
%%\unnumbered% uncomment this for unnumbered level heads

\begin{document}

\title[Article Title]{ Análise de Estatísticas de Jogo e sua Relação com os Resultados no Campeonato Brasileiro }

%%=============================================================%%
%% GivenName	-> \fnm{Joergen W.}
%% Particle	-> \spfx{van der} -> surname prefix
%% FamilyName	-> \sur{Ploeg}
%% Suffix	-> \sfx{IV}
%% \author*[1,2]{\fnm{Joergen W.} \spfx{van der} \sur{Ploeg} 
%%  \sfx{IV}}\email{iauthor@gmail.com}
%%=============================================================%%

\author*[1,2]{\fnm{First} \sur{Caio César Souza Oliveira}}\email{caio.s.oliveira@unesp.br}

\author[2,3]{\fnm{Second} \sur{Igor dos Reis Gomes}}\email{reis.gomes@unesp.br}
\equalcont{Esses autores contribuíram igualmente para este trabalho.}

\affil*[1]{\orgdiv{Bacharelado em Ciência da Computação}, \orgname{UNESP}, \orgaddress{\street{Av. Eng. Luís Edmundo Carrijo Coube, 14-01 - Vargem Limpa}, \city{Bauru}, \postcode{17033-360}, \state{São Paulo}, \country{Brasil}}}

%%==================================%%
%% Sample for unstructured abstract %%
%%==================================%%


\abstract{Este trabalho apresenta uma análise exploratória das partidas do Campeonato Brasileiro no período de 2003 a 2023. A análise utiliza um conjunto de dados público do \textit{Kaggle} e integra informações de partidas, gols, cartões e estatísticas de jogo. A partir disso, foram realizados procedimentos estatísticos para se analisar as informações dos jogos com diversos tipos de testes de hipótese, como teste Qui-quadrado, tamanho de efeito entre outros para assim poder responder a pergunta de que se jogar os jogos em casa altera a probabilidade da vitória.}

%%================================%%
%% Sample for structured abstract %%
%%================================%%

% \abstract{\textbf{Purpose:} The abstract serves both as a general introduction to the topic and as a brief, non-technical summary of the main results and their implications. The abstract must not include subheadings (unless expressly permitted in the journal's Instructions to Authors), equations or citations. As a guide the abstract should not exceed 200 words. Most journals do not set a hard limit however authors are advised to check the author instructions for the journal they are submitting to.
% 
% \textbf{Methods:} The abstract serves both as a general introduction to the topic and as a brief, non-technical summary of the main results and their implications. The abstract must not include subheadings (unless expressly permitted in the journal's Instructions to Authors), equations or citations. As a guide the abstract should not exceed 200 words. Most journals do not set a hard limit however authors are advised to check the author instructions for the journal they are submitting to.
% 
% \textbf{Results:} The abstract serves both as a general introduction to the topic and as a brief, non-technical summary of the main results and their implications. The abstract must not include subheadings (unless expressly permitted in the journal's Instructions to Authors), equations or citations. As a guide the abstract should not exceed 200 words. Most journals do not set a hard limit however authors are advised to check the author instructions for the journal they are submitting to.
% 
% \textbf{Conclusion:} The abstract serves both as a general introduction to the topic and as a brief, non-technical summary of the main results and their implications. The abstract must not include subheadings (unless expressly permitted in the journal's Instructions to Authors), equations or citations. As a guide the abstract should not exceed 200 words. Most journals do not set a hard limit however authors are advised to check the author instructions for the journal they are submitting to.}

\keywords{análise exploratória de dados, teste de hipótese, Campeonato Brasileiro, mando de campo, futebol, ciência de dados}

%%\pacs[JEL Classification]{D8, H51}

%%\pacs[MSC Classification]{35A01, 65L10, 65L12, 65L20, 65L70}

\maketitle

\section{Introdução}\label{sec1}

O futebol é um esporte caracterizado pela interação de diversos fatores que podem influenciar o resultado de uma partida. Embora a vitória ou a derrota seja determinada pelo número de gols marcados, o desempenho das equipes ao longo do jogo pode ser descrito por diferentes estatísticas, como posse de bola, número de finalizações, chutes no alvo, escanteios, passes, faltas e cartões. A análise conjunta dessas informações permite investigar quais características do desempenho de uma equipe estão mais associadas à obtenção de resultados positivos.

Nesse contexto, a Ciência de Dados oferece métodos para explorar grandes conjuntos de informações esportivas e identificar padrões que dificilmente seriam observados apenas pela análise individual das partidas. A utilização de técnicas de análise exploratória de dados permite inicialmente compreender a distribuição das variáveis, identificar possíveis relações entre as estatísticas e os resultados das partidas e observar diferenças no comportamento das equipes como mandantes e visitantes.

Este estudo utiliza um conjunto de dados do Campeonato Brasileiro de Futebol, contendo informações sobre partidas disputadas ao longo de diferentes temporadas, incluindo o resultado dos jogos e diversas estatísticas de desempenho das equipes. A partir desses dados, busca-se analisar a relação entre as características das partidas e o resultado obtido por cada equipe, considerando tanto indicadores ofensivos quanto aspectos relacionados à posse de bola, criação de oportunidades, disciplina e desempenho coletivo.

O objetivo é, portanto, explorar os dados e identificar quais estatísticas apresentam maior associação com a vitória ou derrota das equipes no Campeonato Brasileiro. Para isso, inicialmente é realizada uma análise exploratória dos dados, permitindo compreender suas características e estabelecer uma base para análises posteriores utilizando métodos estatísticos e técnicas de aprendizado de máquina. Dessa forma, o estudo procura transformar os registros das partidas em informações capazes de contribuir para uma melhor compreensão dos fatores relacionados ao desempenho das equipes.

\section{Desenvolvimento}

\subsection{Fonte dos Dados}

O estudo utiliza conjuntos de dados referentes a partidas do Campeonato Brasileiro, organizados de forma a contemplar diferentes aspectos dos jogos. A base \textit{campeonato-brasileiro-full.csv} reúne informações gerais das partidas, incluindo as equipes participantes e seus respectivos resultados. Esse conjunto fornece a estrutura principal para identificar as partidas analisadas e estabelecer a relação entre as equipes e os resultados obtidos. A partir dessas informações, torna-se possível classificar os jogos de acordo com vitória, derrota ou empate. Dessa maneira, a base geral é utilizada como referência para a organização das demais informações empregadas no estudo. Sua utilização permite estabelecer uma visão abrangente dos confrontos analisados ao longo da competição.

Complementarmente, o arquivo \textit{campeonato-brasileiro-estatisticas-full.csv} apresenta informações relacionadas ao desempenho das equipes durante as partidas. Entre as estatísticas consideradas estão aspectos como posse de bola, finalizações, finalizações no alvo, passes, faltas, cartões, impedimentos e escanteios. Esses dados são particularmente relevantes para o objetivo do estudo, pois permitem investigar a relação entre o desempenho apresentado em campo e o resultado final das partidas. A comparação dessas variáveis entre equipes vencedoras e derrotadas possibilita identificar padrões e diferenças no comportamento estatístico dos times. Dessa forma, o conjunto de estatísticas constitui uma das principais fontes para a análise exploratória e para a aplicação dos métodos de Ciência de Dados.

Os arquivos \textit{campeonato-brasileiro-gols.csv} e \textit{campeonato-brasileiro-cartoes.csv} complementam as informações anteriores ao detalhar eventos ocorridos durante as partidas. A base de gols permite analisar os registros de gols considerando informações relacionadas às equipes, aos atletas e aos momentos em que os eventos ocorreram. Já a base de cartões acrescenta informações referentes aos cartões aplicados durante os confrontos, permitindo observar aspectos disciplinares das equipes. A integração dessas diferentes fontes possibilita uma análise mais ampla das partidas, relacionando o resultado final a diferentes características do desempenho esportivo. Assim, os conjuntos de dados fornecem informações complementares que podem ser utilizadas conjuntamente para investigar quais estatísticas apresentam relação com a ocorrência de vitórias e derrotas no Campeonato Brasileiro.

\subsection{Tratamento dos Dados}\label{subsec_tratamento}

O tratamento dos dados foi realizado como etapa inicial da Análise Exploratória de Dados (EDA), com o objetivo de integrar e padronizar as diferentes bases utilizadas no estudo. As fontes apresentavam diferentes níveis de granularidade, sendo algumas organizadas por partida, outras por gol, cartão ou estatística individual de cada equipe. Para possibilitar a comparação entre mandante e visitante, essas informações foram reorganizadas de forma que cada registro representasse uma única partida, contendo os dados correspondentes às duas equipes. Esse procedimento permitiu estabelecer uma estrutura comum entre as diferentes fontes e facilitou a realização das análises posteriores. Após a integração, foram consideradas variáveis relacionadas ao resultado da partida, gols, estatísticas de desempenho e eventos ocorridos durante os jogos. Dessa maneira, a etapa de preparação transformou as bases originais em um conjunto de dados adequado à análise estatística.

Durante a exploração inicial, também foi realizada uma análise da presença de valores ausentes. A identificação desses valores foi importante porque algumas estatísticas não estavam disponíveis para todas as temporadas analisadas. As estatísticas detalhadas de jogo apresentaram dados reais principalmente no período de 2015 a 2023, enquanto as informações de gols e cartões possuem cobertura a partir de períodos diferentes. Dessa forma, os valores que anteriormente apareciam como zero, mas representavam ausência de registro, foram convertidos para valores nulos (`NaN`). Essa alteração evita que partidas sem informação sejam interpretadas incorretamente como partidas nas quais determinado evento não ocorreu. Além disso, o uso dos valores nulos permite que operações estatísticas, como o cálculo de médias, desconsiderem automaticamente as observações sem dados disponíveis. Assim, o tratamento preserva a diferença entre uma ocorrência igual a zero e uma informação que não foi coletada.

Após a padronização e o tratamento dos valores ausentes, o conjunto de dados foi utilizado para a realização da análise exploratória. Foram calculadas estatísticas descritivas relacionadas aos gols do mandante e do visitante, ao total de gols e ao saldo da partida, além da distribuição dos resultados entre vitória do mandante, empate e vitória do visitante. A EDA também possibilitou comparar diferentes características de desempenho entre as equipes, como chutes, chutes no alvo, posse de bola, passes, faltas, impedimentos, escanteios e cartões. Para as estatísticas disponíveis entre 2015 e 2023, foram consideradas 3.419 partidas, enquanto a base geral tratada contempla 8.405 partidas entre 2003 e 2023. Por fim, o conjunto tratado foi exportado para um novo arquivo, chamado \texttt{`campeonato\_brasileiro\_tratado.csv`}, permitindo que os dados preparados fossem usados nas etapas seguintes.

\section{Resultados}\label{sec2}

A análise seguiu duas etapas complementares. Primeiro, uma análise exploratória de dados (EDA) descreveu a magnitude e a estrutura do efeito de mando ao longo das 8.405 partidas do período 2003--2023. Em seguida, quatro testes de hipótese formais, com nível de significância $\alpha = 0{,}05$, avaliaram se esse efeito é estatisticamente distinguível do acaso e investigaram possíveis mecanismos explicativos, usando as 3.419 partidas com estatísticas de jogo disponíveis (2015--2023) sempre que aplicável. Os testes foram implementados em Python, com as bibliotecas \texttt{pandas} e \texttt{scipy.stats}.

\subsection{Distribuição dos resultados}\label{subsec_dist}

Das 8.405 partidas analisadas, o mandante venceu 49,7\% (4.179 partidas), houve empate em 26,4\% (2.221) e o visitante venceu em 23,9\% (2.005) --- uma razão de 2,08 vitórias do mandante para cada vitória do visitante (Figura~\ref{fig_dist}). As estatísticas descritivas de gols reforçam esse padrão: o mandante marca em média 1,54 gol por partida (dp~1,23) contra 1,03 do visitante (dp~1,03), resultando em saldo médio do mandante de 0,51 gol (dp~1,56). O mandante marca ao menos um gol em 79,3\% das partidas e não sofre gols em 36,4\% delas, contra 63,6\% e 20,7\% do visitante, respectivamente --- uma razão de 1,76$\times$ na defesa, superior à razão de 1,25$\times$ observada no ataque.

\begin{table}[h]
\caption{Estatísticas descritivas das variáveis de desfecho (8.405 partidas, 2003--2023)}\label{tab_desc}
\begin{tabular}{@{}lcccc@{}}
\toprule
Variável & Média & Desvio-padrão & Mín. & Máx.\\
\midrule
Gols do mandante   & 1,54 & 1,23 & 0 & 7\\
Gols do visitante  & 1,03 & 1,03 & 0 & 7\\
Total de gols       & 2,57 & 1,64 & 0 & 11\\
Saldo do mandante   & 0,51 & 1,56 & $-$7 & 7\\
\botrule
\end{tabular}
\end{table}

\begin{figure}[h]
\centering
\includegraphics[width=0.7\textwidth]{figuras/distribuicao_de_resultados.png}
\caption{Distribuição dos resultados das partidas do Campeonato Brasileiro (2003--2023).}\label{fig_dist}
\end{figure}

\subsection{Estabilidade da vantagem ao longo do tempo}\label{subsec_tempo}

A vantagem de mando não é um artefato de poucas temporadas: o mandante teve percentual de vitórias superior ao do visitante em todas as 21 temporadas analisadas. A taxa de vitória do mandante variou entre 43,7\% (2020) e 54,7\% (2008), com desvio-padrão de apenas 3,3 pontos percentuais entre temporadas (Figura~\ref{fig_tempo}). Quando a vantagem é medida em gols por partida em vez de vitórias, observa-se uma tendência de queda: a diferença média entre gols do mandante e do visitante caiu de 0,597 gol (2003--2009) para 0,403 gol (2017--2023), uma redução de aproximadamente 33\%, com inclinação linear de $-0{,}0123$ gol por temporada.

\begin{figure}[h]
\centering
\includegraphics[width=0.85\textwidth]{figuras/mando_ao_longo_do_tempo.png}
\caption{Percentual de vitória do mandante, empate e vitória do visitante por temporada.}\label{fig_tempo}
\end{figure}

\subsection{Teste binomial exato: o mando altera a probabilidade de vitória?}\label{subsec_binom}

Para isolar o efeito do mando do acaso, testou-se se, entre as 6.184 partidas decididas (sem empate), o mandante venceria com probabilidade igual a 0,5, como esperado sob uma ``moeda honesta''.

\textbf{H0:} entre partidas decididas, o mandante vence com probabilidade 0,5.\\
\textbf{H1:} a probabilidade de vitória do mandante é diferente de 0,5.

O mandante venceu 4.179 das 6.184 partidas decididas (67,6\%), um desvio de 27,6 desvios-padrão em relação ao valor esperado sob H0 (IC 95\% $[66{,}4\%, 68{,}7\%]$; $p = 1{,}15\times10^{-171}$). H0 é rejeitada com ampla margem: jogar em casa altera de forma inequívoca a probabilidade de vitória.

\subsection{Tamanho de efeito}\label{subsec_efeito}

O teste binomial estabelece que o efeito existe; esta seção o quantifica em duas escalas. A diferença de proporções entre vitória do mandante (49,7\%) e vitória do visitante (23,9\%) é de 25,9 pontos percentuais (IC 95\% $[24{,}5\%, 27{,}3\%]$). Traduzindo para a regra de pontuação do campeonato (3 pontos por vitória, 1 por empate), o mando vale 0,776 ponto por partida (IC 95\% $[0{,}725, 0{,}832]$, obtido por bootstrap com 2.000 reamostragens). Ao longo de uma temporada com 19 jogos em casa, isso equivale a aproximadamente 14,7 pontos adicionais (IC 95\% $[13{,}8, 15{,}8]$) --- uma diferença da ordem de várias posições na tabela final.

\subsection{Teste t pareado: como o mandante joga diferente}\label{subsec_tpareado}

Para investigar o mecanismo por trás da vantagem, comparou-se estatisticamente, partida a partida, o desempenho do mandante e do visitante no mesmo jogo (n = 3.419, 2015--2023), usando teste t de Student pareado. Duas métricas foram escolhidas por captarem aspectos distintos do jogo: chutes (volume ofensivo) e posse de bola (controle de jogo).

\begin{table}[h]
\caption{Comparação mandante vs.\ visitante, teste t pareado (n = 3.419, 2015--2023)}\label{tab_tpareado}
\begin{tabular}{@{}lccccc@{}}
\toprule
Métrica & Mandante & Visitante & Razão M/V & $p$ & Cohen's $d$\\
\midrule
Chutes         & 14,52 & 11,10 & 1,31 & $3{,}05\times10^{-140}$ & 0,452\\
Posse de bola  & 51,72\% & 48,28\% & 1,07 & $2{,}55\times10^{-21}$ & 0,164\\
\botrule
\end{tabular}
\end{table}

Em ambas as métricas H0 (diferença média igual a zero) é rejeitada, mas as magnitudes divergem: o mandante chuta 31\% mais, com efeito de tamanho médio ($d = 0{,}45$), enquanto a posse de bola sobe apenas 7\%, com efeito pequeno ($d = 0{,}16$). Combinado ao fato de que a taxa de conversão de chutes em chutes no alvo é praticamente idêntica entre mandante (26,2\%) e visitante (25,9\%), o padrão sugere que o mandante não joga com mais qualidade técnica, mas ataca com mais volume e iniciativa.

Um padrão semelhante aparece na disciplina: o número médio de faltas cometidas é praticamente idêntico entre mandante (14,85) e visitante (14,84), mas o visitante recebe 11,4\% mais cartões amarelos (2,50 vs.\ 2,24 por jogo) e 42,0\% mais cartões vermelhos (0,150 vs.\ 0,106 por jogo). Normalizando pelo número de faltas, o visitante é advertido a cada 5,94 faltas cometidas, contra 6,62 do mandante.

\begin{table}[h]
\caption{Estatísticas de jogo, mandante vs.\ visitante (média por partida, n = 3.419)}\label{tab_stats}
\begin{tabular}{@{}lccc@{}}
\toprule
Métrica & Mandante & Visitante & Razão M/V\\
\midrule
Chutes no alvo        & 3,81  & 2,88  & 1,32\\
Escanteios            & 5,91  & 4,49  & 1,32\\
Passes                & 431,3 & 406,8 & 1,06\\
Faltas                & 14,85 & 14,84 & 1,00\\
Cartão amarelo/jogo   & 2,24  & 2,50  & 0,90\\
Cartão vermelho/jogo  & 0,106 & 0,150 & 0,70\\
\botrule
\end{tabular}
\end{table}

\subsection{Assimetria em placares elásticos}\label{subsec_placares}

Como verificação adicional e independente dos testes formais, comparou-se a frequência de cada placar com a de seu espelho (por exemplo, 2$\times$1 vs.\ 1$\times$2), o que isola o efeito de mando sem exigir nenhum modelo estatístico. Nenhum par é simétrico, e a razão cresce com a margem de gols (Tabela~\ref{tab_placares}, Figura~\ref{fig_placares}): de 1,86$\times$ em vitórias por 1$\times$0 a 4,66$\times$ em vitórias por 4$\times$0. Esse padrão é consistente com um efeito que se acumula ao longo da partida --- pressão territorial sustentada, retração do adversário após sofrer o primeiro gol --- em vez de um efeito pontual e único.

\begin{table}[h]
\caption{Frequência de placares e seus espelhos (2003--2023)}\label{tab_placares}
\begin{tabular}{@{}lccc@{}}
\toprule
Placar & Mandante vence & Visitante vence & Razão\\
\midrule
1$\times$0 & 1.110 & 598 & 1,86$\times$\\
2$\times$0 & 709   & 308 & 2,30$\times$\\
3$\times$0 & 375   & 122 & 3,07$\times$\\
4$\times$0 & 149   & 32  & 4,66$\times$\\
2$\times$1 & 812   & 512 & 1,59$\times$\\
3$\times$1 & 369   & 149 & 2,48$\times$\\
\botrule
\end{tabular}
\end{table}

\begin{figure}[h]
\centering
\includegraphics[width=0.85\textwidth]{figuras/assimetria_placares.png}
\caption{Frequência de placares a favor do mandante e do visitante, por margem de gols.}\label{fig_placares}
\end{figure}

\subsection{Vantagem por clube}\label{subsec_clube}

Para descartar a hipótese de que o efeito agregado decorre da composição dos confrontos (times fortes jogando mais em casa contra times fracos), comparou-se o aproveitamento de pontos de cada clube em casa e fora, usando o próprio clube como controle. Entre os 32 clubes com ao menos 76 partidas como mandante, o diferencial médio foi de 25,6 pontos percentuais (dp~4,9), variando de 11,3~p.p.\ (Atlético-GO) a 35,4~p.p.\ (Athletico-PR). Nenhum dos 32 clubes analisados apresentou diferencial negativo.

A vantagem de mando também se mostrou praticamente independente da força geral do clube: a correlação entre o aproveitamento médio de pontos do clube e seu diferencial casa--fora foi de $r = -0{,}02$ ($R^2 = 0{,}00$) (Figura~\ref{fig_forca}). Times mais fortes vencem mais tanto em casa quanto fora, mas o acréscimo proporcionado pelo mando é essencialmente o mesmo em todos os quartis de força (variando entre 25,0 e 27,2~p.p.).

\begin{figure}[h]
\centering
\includegraphics[width=0.85\textwidth]{figuras/forca_vs_mando.png}
\caption{Relação entre a força geral do clube (aproveitamento médio) e seu diferencial de mando (casa $-$ fora).}\label{fig_forca}
\end{figure}

\subsection{Qui-quadrado: o efeito do público (portões fechados)}\label{subsec_chi2}

Entre agosto de 2020 e setembro de 2021, partidas do Campeonato Brasileiro foram disputadas sem público por determinação sanitária, oferecendo um quase-experimento para investigar se a torcida é um dos mecanismos da vantagem de mando. Comparou-se a distribuição de resultados em três períodos --- com público, portões fechados e retorno do público --- por meio de um teste qui-quadrado de homogeneidade (140 partidas de transição foram excluídas).

\textbf{H0:} a distribuição dos resultados é a mesma nos três períodos.\\
\textbf{H1:} ao menos um período difere.

\begin{table}[h]
\caption{Taxa de vitória do mandante por período de acesso do público}\label{tab_chi2}
\begin{tabular}{@{}lccc@{}}
\toprule
Período & $n$ & Vitória do mandante & IC 95\%\\
\midrule
Com público (até mar/2020)            & 6.885 & 50,7\% & $[49{,}5\%, 51{,}8\%]$\\
Portões fechados (ago/2020--set/2021) & 594   & 42,4\% & $[38{,}4\%, 46{,}5\%]$\\
Público de volta (dez/2021 em diante) & 786   & 45,9\% & $[42{,}4\%, 49{,}5\%]$\\
\botrule
\end{tabular}
\end{table}

O teste rejeita H0 ($\chi^2 = 20{,}98$, gl~$=4$, $p = 3{,}20\times10^{-4}$; Figura~\ref{fig_portoes}). A queda de 8,3 pontos percentuais entre o período com público e o período sem público é maior que a incerteza das duas estimativas, e é compatível com a hipótese de que a torcida seja um componente relevante da vantagem de mando. No entanto, o efeito não desaparece: mesmo sem público, o mandante venceu 42,4\% contra 27,6\% do visitante, restando uma vantagem substancial atribuível a mecanismos que não dependem de torcida, como familiaridade com o campo e desgaste da viagem do visitante. O intervalo do período de retorno do público (45,9\%) se sobrepõe tanto ao período anterior quanto ao posterior, de modo que a recuperação da vantagem é sugestiva, mas não conclusiva.

\begin{figure}[h]
\centering
\includegraphics[width=0.7\textwidth]{figuras/portoes_fechados.png}
\caption{Taxa de vitória do mandante nos três períodos de acesso do público, com IC 95\%.}\label{fig_portoes}
\end{figure}

\subsection{Síntese dos testes}\label{subsec_sintese}

\begin{table}[h]
\caption{Síntese dos quatro testes de hipótese realizados}\label{tab_sintese}
\begin{tabular}{@{}p{2.3cm}p{3.5cm}p{4.3cm}c@{}}
\toprule
Teste & H0 & Resultado & Decisão\\
\midrule
Binomial exato & Mandante vence 50\% das decididas & 67,6\%, IC95 $[66{,}4\%, 68{,}7\%]$; $p \approx 1{,}1\times10^{-171}$ & Rejeita H0\\
$t$ pareado & Médias iguais em chutes e posse & $p<0{,}001$ nas duas; $d=0{,}45$ e $0{,}16$ & Rejeita H0\\
Qui-quadrado & Resultados iguais nos três períodos & $50{,}7\% \to 42{,}4\% \to 45{,}9\%$; $p=3{,}2\times10^{-4}$ & Rejeita H0\\
\botrule
\end{tabular}
\end{table}

\section{Discussão}\label{sec12}

Os resultados obtidos convergem para um quadro consistente sobre a vantagem de mando no Campeonato Brasileiro. Em primeiro lugar, o efeito é robusto: está presente em todas as 21 temporadas analisadas e em todos os 32 clubes com amostra suficiente, e é estatisticamente indistinguível do acaso apenas sob a hipótese nula de ausência de vantagem --- hipótese amplamente rejeitada pelo teste binomial exato.

Em segundo lugar, os dados favorecem uma explicação baseada em volume e iniciativa territorial, e não em superioridade técnica. O mandante chuta cerca de 31\% mais que o visitante, mas converte esses chutes em gols na mesma taxa (26,2\% contra 25,9\%), e sua posse de bola é apenas marginalmente maior. Esse padrão é mais compatível com pressão territorial sustentada do que com um time que ``joga melhor'' em casa.

Em terceiro lugar, a assimetria disciplinar --- faltas praticamente idênticas, mas cartões amarelos e vermelhos sistematicamente mais frequentes contra o visitante --- é consistente com viés de arbitragem sob pressão de torcida, hipótese já discutida na literatura de economia do esporte, embora explicações alternativas (diferenças qualitativas no tipo de falta cometida por quem defende mais, ou maior contestação de quem está perdendo, já que o visitante perde mais) não possam ser descartadas com os dados disponíveis.

O quase-experimento dos portões fechados oferece a evidência mais direta sobre o papel da torcida: a vantagem de mando caiu de forma estatisticamente significativa na ausência de público, mas não desapareceu. Isso sugere que a vantagem de mando é multicausal, combinando um componente ligado à torcida com componentes estruturais (familiaridade com o campo, desgaste de viagem do visitante) que persistem independentemente do público.

A ausência de correlação entre a força do clube e o tamanho de sua vantagem de mando é outro resultado notável: contraria a intuição de que clubes grandes, com estádios maiores e torcidas mais numerosas, extrairiam mais vantagem do mando. Os dados indicam, ao contrário, que o mando funciona como um efeito aditivo e aproximadamente uniforme, somando cerca de 25 pontos percentuais ao aproveitamento de qualquer clube, independentemente de sua força.

\textbf{Limitações.} Os intervalos de confiança reportados pressupõem partidas independentes; há dependência residual não modelada (adversários de forças diferentes, autocorrelação de desempenho ao longo da temporada), de modo que os intervalos reais são provavelmente um pouco mais largos que os apresentados. As estatísticas de jogo (chutes, posse, faltas) só têm cobertura real a partir de 2015, restringindo a Seção~\ref{subsec_tpareado} a nove das 21 temporadas. O período de portões fechados coincidiu com uma temporada atípica em múltiplos eixos (calendário condensado, ausência de pré-temporada, desfalques por surtos de COVID-19), o que impede isolar cirurgicamente o efeito da ausência de torcida do efeito de uma temporada anômala. Por fim, nenhuma das análises estabelece causalidade em sentido estrito; o desenho de portões fechados é o mais próximo de um experimento controlado que os dados permitem, mas ainda carrega as ressalvas anteriores. A série é restrita ao Campeonato Brasileiro, e a generalização para outras ligas exigiria replicação.

\section{Conclusão}\label{sec13}

Este trabalho buscou responder se jogar em casa altera a probabilidade de vitória no Campeonato Brasileiro e, em caso afirmativo, por quanto e por quê. A resposta à primeira pergunta é inequívoca: entre partidas decididas, o mandante vence 67,6\% delas (IC 95\% $[66{,}4\%, 68{,}7\%]$), muito acima dos 50\% esperados sob ausência de vantagem ($p \approx 1{,}15\times10^{-171}$), efeito positivo em todas as 21 temporadas e em todos os 32 clubes analisados.

Quanto à magnitude, o mando acrescenta cerca de 26 pontos percentuais à probabilidade de vitória e vale aproximadamente 0,78 ponto por partida, equivalente a cerca de 14,7 pontos ao longo de uma temporada com 19 jogos em casa --- uma diferença capaz de mover um clube várias posições na tabela final.

Quanto ao mecanismo, os dados indicam que o mandante ataca com mais volume (31\% mais chutes) sem ganho de eficiência, e é beneficiado por uma assimetria disciplinar na arbitragem (42\% mais cartões vermelhos contra o visitante). O quase-experimento dos portões fechados mostra que a torcida é um componente relevante --- a vantagem caiu de 50,7\% para 42,4\% na ausência de público --- mas não o único, já que uma vantagem substancial permanece mesmo sem torcida nos estádios. A vantagem, além disso, mostrou-se estatisticamente estável ao longo do tempo em termos relativos, ainda que a diferença de gols tenha diminuído cerca de um terço entre o início e o fim da série, e mostrou-se uniforme entre clubes de diferentes forças, sem evidência de que times maiores se beneficiem mais do fator casa.

Como direções futuras, seria útil (i) modelar a dependência entre partidas de uma mesma temporada para refinar os intervalos de confiança; (ii) explorar dados de distância de viagem e altitude para quantificar separadamente o componente logístico da vantagem de mando; e (iii) replicar as análises em outras ligas sul-americanas para avaliar até que ponto os padrões aqui descritos são específicos do futebol brasileiro.


\backmatter

\end{document}